\begin{abstract}
Few-shot adaptation of biomedical vision-language models depends on which examples are drawn.
Prompt-Anchored Residual Adaptation (PARA) addresses this by keeping the frozen prompt prediction as an anchor
and moving towards a prediction trained on the support set by a step computed from the geometry of the support embeddings.
We reimplement PARA on BiomedCLIP and study the step on six MedMNIST datasets, three used for design and three held out,
with repeated support draws, balanced accuracy and worst-case (CVaR20) evaluation.
We find that the geometric signals barely depend on the class in BiomedCLIP's frozen space,
so the step reduces to the same function of the number of shots on every dataset ($0.33$, $0.60$ and $0.80$ at 1, 4 and 16 shots),
predictable from two constants of the pretrained model; a schedule in the shot count alone does better.
On binary tasks, standardizing each branch's scores turns the blend into a switch: PARA returns the zero-shot prediction below a step of $0.5$
and the trained branch's prediction above it, on every image.
The best step instead follows how reliable the anchor is relative to the trained branch;
a rule based on the anchor's reliability alone, fixed before the held-out runs, improved on the design datasets and failed on the held-out ones.
Blending log-probabilities instead of standardized scores, with no new parameters, raises balanced accuracy by $4.86$ points on multi-class datasets
but lowers the worst case where the anchor is good.
We conclude with recommendations for evaluating adaptive blending: check that adaptive signals vary across support draws,
compare against shot-only and single-branch baselines, include binary tasks, and test on held-out data.
\end{abstract}
